\documentclass[12pt,notitlepage]{report}
% The shared notes preamble lives outside this repo.
\input{/Users/aryankulkarni/Documents/y3-notes/preamble.tex}

\title{Week 1: exploring the data and a first linear model}
\author{Aryan Kulkarni}

\begin{document}
\maketitle

% Run the R blocks in order with <space>rr. R's working directory is this
% directory, so figures/ and output/ mean the same to R and to LaTeX, and the
% data sits two levels up. The out() helper saves printed output to output/,
% which the text blocks below include.

\section{Data and variable types}

\begin{minted}{r}
library(tidyverse)

matches <- read_csv("../../ashford_matches.csv") |>
  mutate(home_away = factor(home_away))

players <- read_csv("../../league_players.csv") |>
  mutate(across(c(team, position, injured), factor))

out <- function(x, name) {
  writeLines(capture.output(x), sprintf("output/%s.txt", name))
}
\end{minted}

\begin{minted}{r}
out(glimpse(matches), "glimpse-matches")
out(summary(matches), "summary-matches")
\end{minted}

\inputminted{text}{output/glimpse-matches.txt}
\inputminted{text}{output/summary-matches.txt}

\begin{minted}{r}
out(glimpse(players), "glimpse-players")
out(summary(players), "summary-players")
\end{minted}

\inputminted{text}{output/glimpse-players.txt}
\inputminted{text}{output/summary-players.txt}

\paragraph{Classification.}
\mintinline{r}|read_csv| only decides how a column is stored (numbers or
text), so the types below are our own reading of what each variable means.

\begin{center}
  \begin{tabular}{lll}
    Data set & Type & Variables \\
    \toprule
    matches & categorical      & \mintinline{r}|home_away| \\
    matches & discrete         & \mintinline{r}|goals_scored|, \mintinline{r}|goals_conceded| \\
    matches & discrete ordinal & \mintinline{r}|opponent_rank| \\
    players & categorical      & \mintinline{r}|team|, \mintinline{r}|position| \\
    players & binary           & \mintinline{r}|injured| \\
    players & discrete         & \mintinline{r}|prev_injuries| \\
    players & discrete ordinal & \mintinline{r}|training_load| \\
    players & continuous       & \mintinline{r}|bmi|, \mintinline{r}|age|, \mintinline{r}|minutes_played| \\
  \end{tabular}
\end{center}

\section{Plots before modelling}

\begin{minted}{r}
p <- ggplot(players, aes(x = minutes_played)) +
  geom_histogram(bins = 30)
ggsave("figures/hist-minutes.pdf", p, width = 6, height = 4)

p <- ggplot(players, aes(x = bmi)) +
  geom_histogram(bins = 30)
ggsave("figures/hist-bmi.pdf", p, width = 6, height = 4)
\end{minted}

\incfig[0.8]{hist-minutes}

\incfig[0.8]{hist-bmi}

There is no heavy skew or pile-up at the edges of the histograms, so
approximately normal residuals are plausible.

\begin{minted}{r}
p <- ggplot(players, aes(x = position, y = minutes_played)) +
  geom_boxplot()
ggsave("figures/box-minutes-position.pdf", p, width = 6, height = 4)
\end{minted}

\incfig[0.8]{box-minutes-position}

The data is fairly homogeneous within each position.

\begin{minted}{r}
p <- ggplot(players, aes(x = age, y = minutes_played, colour = position)) +
  geom_point(alpha = 0.5) +
  geom_smooth(method = "lm", se = FALSE)
ggsave("figures/scatter-minutes-age.pdf", p, width = 6, height = 4)
\end{minted}

\incfig[0.8]{scatter-minutes-age}

There is no bending in the age scatter, so a linear \mintinline{r}|age| term
makes sense.

\begin{minted}{r}
p <- ggplot(players, aes(x = factor(training_load), y = minutes_played)) +
  geom_boxplot()
ggsave("figures/box-minutes-load.pdf", p, width = 6, height = 4)
\end{minted}

\incfig[0.8]{box-minutes-load}

The \mintinline{r}|training_load| boxes rise evenly, so a numeric term is fine.

\begin{minted}{r}
p <- ggplot(players, aes(x = position, y = bmi)) +
  geom_boxplot()
ggsave("figures/box-bmi-position.pdf", p, width = 6, height = 4)
\end{minted}

\incfig[0.8]{box-bmi-position}

\section{Linear model for minutes played}

Injury is left out of the model because it directly reduces
\mintinline{r}|minutes_played|, so it is a consequence of the outcome.

\begin{minted}{r}
fit_min <- lm(
  minutes_played ~ position + age + training_load + prev_injuries,
  data = players
)
out(summary(fit_min), "summary-fit-min")
\end{minted}

\inputminted{text}{output/summary-fit-min.txt}

These explanatory variables explain about 27\% of the variation in minutes
played. The residual standard error is 605 minutes, so a typical prediction
misses by about 6.7 full matches. \mintinline{r}|position| and
\mintinline{r}|training_load| are very significant predictors of playing
minutes.

\begin{minted}{r}
pdf("figures/diag-min.pdf", width = 7, height = 7)
par(mfrow = c(2, 2))
plot(fit_min)
dev.off()
\end{minted}

\incfig[0.9]{diag-min}

\end{document}
